\documentclass[11pt]{article}
\usepackage{amsmath,amssymb,amsthm,hyperref}
\newtheorem{lemma}{Lemma}
\newtheorem{proposition}{Proposition}
\title{Nonuniform dyadic coordinates: spectrum gains and readout losses}
\author{Research checkpoint 006}
\date{31 October 2026}
\begin{document}
\maketitle

\paragraph{Scope.}
The target remains general-prime-field ECDLP in expected bit time
$2^{\widetilde O(\sqrt b)}$. We test geometric coordinates, repair their
local integrality, and show that a particular odd-node reconstruction
and unrestricted mode-enumeration route still has exponentially many
feasible candidates. This is not a lower bound for all nonuniform grids,
all support-aware contractors, or ECDLP. No improved complete bound is
claimed. The ordinary geometric set also has unequal parity populations;
this checkpoint does not assert a balanced packing theorem for it.

\section{Definitions and a uniform-scaling invariant}
Let $X$ be a finite integer coordinate set, $D=X-X$, and let a finite
integer grid $\Gamma$ contain $D$. Set
\[
 C_\Gamma(T)=\prod_{s\in\Gamma}(T-s),\qquad
 W_{xy}(T)=\frac{C_\Gamma(T)}{C_\Gamma'(0)(T+x-y)}.
\]
The apparent denominator cancels, and $W(0)=I$. Write
\[
 V_\Gamma(t)=v_2(C_\Gamma'(t))
 =\sum_{s\in\Gamma\setminus\{t\}}v_2(t-s).
\]
At a node $\tau=y-x$, the nonzero constant entry has valuation
$V_\Gamma(\tau)-V_\Gamma(0)$. Changing the scale in
$T=\tau+2^\ell Z$ cannot alter this constant.

For any nonzero integer $c$, simultaneously replacing $X,\Gamma,T$
by $cX,c\Gamma,cT$ leaves every $W$ entry unchanged. Indeed the factors
$c^{|\Gamma|}$, $c^{|\Gamma|-1}$ and $c$ cancel. Translating every
coordinate also leaves all differences unchanged. Thus uniform affine
spacing changes alone cannot improve the spectrum or precision tradeoff.

\section{The geometric grid fails local integrality}
For $n\ge3$ put $m=n-2$ and
\[
 X_n=\{0,1,2,4,\ldots,2^m\},\qquad D_n=X_n-X_n.
\]
All grids here are sets of \emph{distinct} differences, not multisets.
The number of positive differences with valuation $i$ is
\[
 c_i=m-i\quad(0\le i<m),\qquad c_m=1.
\]
Consequently
\begin{equation}
 |D_n|=m^2+m+3,\qquad
 V_{D_n}(0)=\frac{m^3+5m}{3}.
 \label{eq:origin}
\end{equation}
The odd positive differences are $2^j-1$, $1\le j\le m$.
Summing their differences from $1$, including their negatives, gives
\begin{equation}
 V_{D_n}(1)=\frac{n^2-n-4}{2}.
 \label{eq:one}
\end{equation}
The difference between \eqref{eq:one} and \eqref{eq:origin} is negative
for every $n\ge3$, with cubic magnitude. For instance $n=8$ gives
$26-82=-56$. Thus the naive geometric deformation is not locally
integral even at the constant term of $W_{0,1}(1+2^\ell Z)$.

\section{A polynomial-size repair, charged rather than dismissed}
For a shift exponent $E$ greater than every old pair-difference
valuation, define
\[
 \Gamma_E=D_n\cup\{d+2^E:d\in D_n\setminus\{0\}\}.
\]
There is no added partner of zero. The grid has $2|D_n|-1=O(n^2)$ nodes.
With $V=V_{D_n}$ and $U=V_{\Gamma_E}$, direct valuation identities give
\begin{equation}
 U(0)=2V(0),\qquad
 U(d)=U(d+2^E)=E+2V(d)-v_2(d)\quad(d\ne0).
 \label{eq:repair}
\end{equation}
Thus constant-term integrality is equivalent to
$E\ge2V(0)-2V(d)+v_2(d)$ for every $d\ne0$.

\begin{lemma}
The smallest $E$ satisfying both strict old-point separation and
the constant-term conditions in this paired construction is
\[
 E=\max\{n,\,2V(0)-6n+16\}.
\]
For $n\ge4$ the second term attains the maximum. In particular
$E=\Theta(n^3)$.
\end{lemma}
\begin{proof}
At valuation level zero the minimum of $V(d)$ is $3m-2$, attained
at $d=\pm3$ if $m\ge2$, and at $d=\pm1$ if $m=1$.
For $j\ge2$, the value on $2^j-1$ is
$jm-j(j-1)/2+m-1$, which is minimized at $j=2$;
the value on $1$ is $(m^2+3m-2)/2$ and is no smaller.

For a nonzero node of valuation $i$, let
\[
 B_i=2\sum_{t=0}^m\min(i,t)c_t.
\]
The minimum derivative valuation on that level is
$B_i+3c_i-2$. This follows by scaling the odd points in that
valuation class; other levels contribute exactly $B_i$.
These minima increase by at least one at each step $i\to i+1$,
and the final step increases by two. Indeed
$B_{i+1}-B_i=2\sum_{t\ge i+1}c_t$; until the last step the odd-class
minimum decreases by three, and at the last step it is unchanged.
The $m=1$ case gives an increase of two directly.
Hence the minimum of $2V(d)-v_2(d)$ is $2(3m-2)=6n-16$.
The largest old pair-difference valuation is $m+1=n-1$, attained
between $2^m$ and $-2^m$. Strict separation requires $E\ge n$.
Combining the two requirements proves the formula.
\end{proof}

Choose $\ell=E$. All differences of distinct repaired nodes have
valuation at most $E$. For $\tau\ne\delta\in D_n$, the local entry
$W_\delta(\tau+2^\ell Z)$ has a factorization with leading factor of
valuation
\[
 \ell+U(\tau)-U(0)-v_2(\tau-\delta)
\]
times $Z$ and integral factors $1+2^\ell Z/(\tau-s)$.
For $\tau=\delta$, its constant prefactor has valuation
$U(\tau)-U(0)\ge0$ and the remaining factors are integral.
Every local polynomial is therefore integral. The repair has polynomial
node count and polynomial coordinate bit length; it is not rejected
merely because some coordinates are numerically large.
For $n\ge4$, this common local scale is minimal for the repaired grid:
at $\tau=3$ one has $U(3)=U(0)$, while the unique paired root at distance
$2^E$ makes the linear coefficient of $W_3(3+2^\ell Z)$ have valuation
$\ell-E$. Every other root difference has smaller valuation, so it
cannot cancel that leading denominator contribution.

For odd nodes set
\[
 \delta_\tau=\ell+U(\tau)-U(0).
\]
The minimum is $E$ for $n\ge4$ (it is $4$ for $n=3$).
The even coordinate block is
$B=(0,2,4,\ldots,2^m)$ of size $n-1$.
At any odd $\tau$, all Cauchy denominators $\tau+x-y$ in this block
are units. Its normalized diagonal increments are exactly
\begin{equation}
 \Delta_k=2\sum_{j<k}v_2(B_k-B_j)=k(k+1),
 \quad 0\le k\le n-2.
 \label{eq:quadratic}
\end{equation}
The associated mode weights are $\delta_\tau+k(k+1)$.
Both normalized triangular factors are integral and have unit diagonal.
Their off-diagonal valuations are $i-k$ for $k>0$, rather than all zero;
this distinction was caught by the independent elimination check.

\section{A realizable readout witness defeats unrestricted enumeration}
Use $M/n$ full packets, $M$ divisible by $4n$, with each old row and
column variable pinned to one. The original contraction is one and its
deformation is
\begin{equation}
 F(T)=\left(\prod_{x\in X_n}\sum_{y\in X_n}W_{xy}(T)\right)^{M/n}.
 \label{eq:realizable}
\end{equation}
This is a valid finite tensor-network deformation, not just an arbitrary
polynomial. It has degree at most $M(|\Gamma_E|-1)$.

Every odd difference of two coordinates in $X_n$ involves its only odd
coordinate, $1$. Hence at most two rows can contain an old difference
equal to a given odd node $d\in D_n$. The same two rows are the only
ones with a difference extremely close to the repaired node $d+2^E$.
For each of the other at least $n-2$ rows and all its columns,
$\tau\ne y-x$ and $v_2(\tau+x-y)\le n$.
The local row form has zero constant term and is coefficientwise
divisible by $2^{\delta_\tau-n}$. The remaining rows are integral.
Thus at every odd repaired node,
\begin{equation}
 \operatorname{ord}_Z F(\tau+2^\ell Z)\ge r:=\frac{M(n-2)}n,
 \qquad
 v_2(\text{every coefficient})\ge
 b_\tau:=\frac{M(n-2)}n(\delta_\tau-n).
 \label{eq:witness}
\end{equation}

Consider reconstruction from \emph{only} odd-node jets, with jet length
$q_\tau$ and precision $P_\tau$. If every node has either
$q_\tau\le r$ or $P_\tau\le b_\tau$, every supplied residue is zero.
Those data cannot distinguish \eqref{eq:realizable}, whose value at zero
is one, from the zero deformation. Some required node therefore has
both $q_\tau>r$ and $P_\tau>b_\tau$.

For $n\ge7$ one has
\[
 \delta_\tau\ge E\ge\frac{n(n-1)(n-2)}3.
\]
For completeness, writing $m=n-2=5+t$, three times the difference
between $E$ and the displayed fraction is
$t^3+12t^2+35t+12>0$.
It follows that
\begin{equation}
 b_\tau>
 \frac M4\bigl(\delta_\tau+(n-1)(n-2)\bigr).
 \label{eq:budget}
\end{equation}
Indeed, after cancelling $M/n$, the needed positivity is
$(3n-8)\delta_\tau-n(n-2)(n+3)>0$, which follows from the preceding
bound and $3n^2-14n-1>0$.

There are $M(n-1)/n$ modes in these even blocks. Choose any fixed $M/4$
of them. Every binary exponent vector supported on that chosen set has
weight below $P_\tau$ by \eqref{eq:budget}, and total degree at most
$M/4<q_\tau$. It also satisfies every per-packet degree cap.
Therefore unrestricted feasible-mode enumeration at that required node
includes at least $2^{M/4}$ candidates.

\paragraph{What this excludes.}
The geometric coordinates, this minimal paired-node repair, and a
reconstruction/contractor using only odd-node residues with unrestricted
candidate enumeration do not obtain a sublinear exponent in $M$ from
the quadratic spectrum. The claim is about that concrete method and
its feasible candidate universe. It does not assert that every candidate
contributes, exclude support-aware pruning, forbid even-node schedules,
or rule out direct evaluation using extra tensor structure. The all-one
network is itself easy to contract; its role is a necessary readout
precision and jet-length test, not a hard instance for every algorithm.

\section{Evidence and next branch}
The independent standard-library script checks sixteen coordinate sizes
$n=3,\ldots,18$, 5024 derivative valuations, direct derivative products,
sixteen exact Cauchy eliminations, twenty-two expanded realizable packet
jets, and 54 uniform-dilation identities. Proofs above, rather than this
finite sampling, establish the all-$n$ statements.

A pilot also tests balanced low-bit cluster coordinates
\[
 X_{u,n}=\{2^u x+r:x\in X_n,\ 0\le r<2^u\}.
\]
These have equal occupancy of all $u$-bit orientations. In fifteen finite
cases, central nodes can be integral even when other retained nodes fail;
for $n=8$ and $u=1,2,3$, the worst retained constant valuation is $-16$.
This does not settle that family. The next test should select admissible
nodes explicitly and jointly charge the larger difference grid, Hermite
multiplicities, local scale, mode weights and packet-entry cost. It must
not inherit the old precision formula without a new derivation.
\end{document}
